\documentclass[11pt]{article}
\usepackage{amsmath,amssymb,amsthm,geometry,xcolor,hyperref}
\hypersetup{colorlinks=true,allcolors=blue!45!black}
\geometry{margin=1in}
\newtheorem{theorem}{Theorem}
\title{A finite-semigroup proof of the permutation-like matrix spectral assertion}
\author{Myeongjin Kang}
\date{October 25, 2026}
\begin{document}
\maketitle
\begin{abstract}
A square matrix with one entry equal to one in each column and all other
entries zero represents a map on a finite set. Finiteness makes two of its
powers equal, proving that its nonzero eigenvalues are roots of unity.
We give the complete argument and a Lean certificate, including the
zero-eigenvalue characterization. This answers the assertion recorded as
MathDB362146 and Conjecture 5.2 of Lippold's paper, without a claim of
worldwide novelty or priority.
\end{abstract}
\begin{theorem}
Let $K$ be a field and $I$ a finite index set. Let $A\in K^{I\times I}$
have exactly one entry equal to one in each column, with all other entries
zero. Every eigenvalue $\mu\in K$ is zero or satisfies $\mu^k=1$ for some
integer $k>0$. Moreover, zero is an eigenvalue if and only if $A$ is not a
permutation matrix.
\end{theorem}
\begin{proof}
For each column $j$, let $f(j)$ be its unique unit-entry row. Then
$A=T_f$, where $(T_f)_{ij}=1$ if $i=f(j)$ and zero otherwise.
Multiplication gives $T_fT_g=T_{f\circ g}$, so $T_f^n=T_{f^n}$.
There are finitely many functions $I\to I$. Hence $f^p=f^q$ for some
integers $0\le p<q$, and $A^p=A^q$.

For an eigenvector $v\ne0$ with $Av=\mu v$, induction gives
$A^nv=\mu^nv$. Thus $\mu^pv=\mu^qv$. A nonzero coordinate of $v$
implies $\mu^p=\mu^q$. If $\mu\ne0$, division gives
$\mu^{q-p}=1$, proving the first assertion.

If $f$ is not injective, choose $i\ne j$ with $f(i)=f(j)$. Then
$T_f(e_i-e_j)=e_{f(i)}-e_{f(j)}=0$, while $e_i-e_j\ne0$.
Thus zero is an eigenvalue. If $f$ is injective, finiteness makes it
bijective, and $T_{f^{-1}}T_f=I$ rules out a nonzero kernel vector.
Finally, $f$ is injective exactly when each row of $T_f$, as well as each
column, has one unit entry; this is the permutation-matrix condition.
\end{proof}
\paragraph{Relationship to the source.}
Lippold \cite{lippold} defines these matrices in Definition 2.1 and proves
their multiplicative closure in Proposition 2.1. Corollary 4.5 gives the
spectral assertion for dimensions two and three, Proposition 4.2 treats
pre-row matrices, and Proposition 4.4 and Corollary 4.7 give the
zero-eigenvalue and invertibility statements. Those are credited antecedents.
The finite-set argument above removes the dimension and pre-row restrictions.
The paper itself notes after Conjecture 5.2 that the zero-eigenvalue part is
already settled. The identity is a permutation matrix, so a separate
exclusion of the identity in the conjecture is redundant.

\paragraph{Formal coverage and limitations.}
\texttt{PermutationSpectrum.lean} proves the full theorem for every field and
finite index type. It constructs the finite-map representation, proves
multiplication and powers, obtains repeated powers by the pigeonhole principle,
and constructs the zero eigenvector or inverse explicitly. The final statement
uses mathlib's \texttt{Module.End.HasEigenvalue}; a separately proved equivalence
connects it with the nonzero-vector equation. The declaration
\texttt{exact\_mathdb\_conclusion} also states the redundant identity exclusion.
No classification of Jordan blocks or algebraic multiplicities is claimed.

The build uses Lean 4.34.1 and mathlib revision
\begin{center}\small\texttt{d13f23b723b8a846827a245b89c10fc7d3f11612}.\end{center}
It prints every local type and
its kernel axioms, and permits only standard axioms. The compiler and
precompiled dependency objects remain trusted. A bounded source search found
no separate general solution; this is not an exhaustive literature review.
Original prose is CC BY 4.0 and original code MIT. No external review or
submission is asserted.
\begin{thebibliography}{9}
\bibitem{lippold} S. R. Lippold, \emph{Permutation-Like Matrices},
arXiv:2403.02478v1, 2024.
\url{https://arxiv.org/html/2403.02478v1}.
\bibitem{mathdb} MathDB362146, \emph{The eigenvalue conjecture for
permutation-like matrices}, retrieved October 25, 2026.
\url{https://mathdb.com/p/362146}.
\end{thebibliography}
\end{document}
